% ch40_localization.tex — Volume IV Chapter 4
% Retrofitted with full Vol I-III pedagogy: cycle stages, etymology, dialogs.

\chapter{Localization at Weak Equivalences}

\begin{overviewbox}
Many constructions of modern category theory follow a common recipe: take
a category $\mathcal{C}$ of ``concrete'' objects, pick out a class
$\mathcal{W}$ of morphisms that ``ought to be isomorphisms even though
they aren't,'' and form a new category $\mathcal{C}[\mathcal{W}^{-1}]$
in which every arrow of $\mathcal{W}$ has been turned into a genuine
isomorphism. The result is the \term{localization} of $\mathcal{C}$
at $\mathcal{W}$.

You have met three motivating special cases. The \emph{derived category}
$\mathcal{D}(\mathcal{A})$ of an abelian category (Chapter~54), where
$\mathcal{W} = $ quasi-isomorphisms; the \emph{homotopy category}
$\mathrm{Ho}(\mathbf{Top})$ of topological spaces, where $\mathcal{W} = $
weak homotopy equivalences; and Gabriel--Zisman's original construction
of ring localizations, where $\mathcal{W}$ is a multiplicatively-closed
subset.

This chapter develops the general theory in three layers. First, the
\term{universal property}: $\mathcal{C}[\mathcal{W}^{-1}]$ is universal
among categories receiving a $\mathcal{W}$-inverting functor from
$\mathcal{C}$. Second, the \term{Gabriel--Zisman construction}: morphisms
are \emph{zigzags} of arrows in $\mathcal{C}$ with formal inverses
inserted. Third, the \term{calculus of fractions}: conditions under
which every zigzag reduces to a single span or cospan. The chapter
closes with the three motivating examples, in preparation for the
model-categorical machinery of Chapter~41 that makes localizations
tractable.
\end{overviewbox}


% ═══════════════════════════════════════════════════════════════════════════
\section{Categories with Weak Equivalences}
\label{sec:weak-equivalences}
% ═══════════════════════════════════════════════════════════════════════════

\subsection{Informally}

A \emph{weak equivalence} is a morphism that ought to be considered an
isomorphism, even though it isn't one in $\mathcal{C}$. Common patterns:
\begin{itemize}
  \item In $\mathbf{Top}$, a continuous map $f : X \to Y$ that induces
        an isomorphism on $\pi_n$ for all $n$ — without being a
        homeomorphism.
  \item In $\mathbf{Ch}(\mathcal{A})$ (chain complexes), a chain map
        that induces an isomorphism on homology — without being
        invertible at the chain level.
  \item In $\mathbf{sSet}$, a map of simplicial sets that becomes a
        homotopy equivalence after geometric realization — without
        having an inverse in $\mathbf{sSet}$.
\end{itemize}
In each case the original category retains too much rigidity. Localizing
forces the weak equivalences to be invertible, recovering the homotopical
content of the situation.

\begin{nameorigin}
The term \term{localization} is borrowed from commutative algebra and
algebraic geometry. To \emph{localize} a ring $R$ at a multiplicative
set $S$ is to invert the elements of $S$ — making them units — producing
$S^{-1} R$. ``Local'' because the result captures the behavior of $R$
near (= local to) the prime ideal $\{r : r \notin S\}$. The categorical
\emph{localization} $\mathcal{C}[\mathcal{W}^{-1}]$ is exactly the same
operation lifted up a level: invert a class of arrows by formally
adjoining inverses. The terminology was carried into homotopy theory
by Bousfield, Quillen, and others in the 1960s--70s.
\end{nameorigin}

\begin{nameorigin}
\term{Weak equivalence} contrasts with the stricter notions of
\emph{isomorphism} or \emph{equivalence of categories}. ``Weak'' here
means ``up to homotopy or some other higher equivalence relation,''
versus ``on the nose.'' The terminology is from algebraic topology:
a \emph{weak homotopy equivalence} of spaces is one inducing isos on
all $\pi_n$ — weaker than an honest homotopy equivalence (Whitehead's
theorem makes the two coincide for CW complexes, but not in general).
\end{nameorigin}

\subsection{Expanding}

Why bother with weak equivalences at all, when isomorphism already
captures the strictest sense of ``same''? Two reasons.

\paragraph{Pragmatic.} In many natural categories, isomorphism is too
strict. Two topological spaces of ``the same homotopy type'' are rarely
homeomorphic; two chain complexes with ``the same homology'' are
rarely chain-isomorphic. The natural notion of sameness is
\emph{homotopy} or \emph{quasi-iso}, not honest isomorphism. Localizing
puts that natural notion on equal footing with strict isomorphism by
forcing the formal collapse.

\paragraph{Structural.} The localization functor $Q : \mathcal{C} \to
\mathcal{C}[\mathcal{W}^{-1}]$ has a universal property:
$\mathcal{W}$-inverting functors out of $\mathcal{C}$ correspond to
ordinary functors out of $\mathcal{C}[\mathcal{W}^{-1}]$. This makes
$\mathcal{C}[\mathcal{W}^{-1}]$ the ``universal home for
$\mathcal{W}$-invariant functors,'' and the construction itself a left
adjoint at the $2$-categorical level.

\begin{dialog}
\textbf{Reader:} Why not just quotient $\mathcal{C}$ by the equivalence
relation ``connected by a zigzag of weak equivalences''?

\textbf{Author:} You can, and that's exactly what the localization does
at the level of objects. But the localization also has to manage
\emph{morphisms}: a morphism $X \to Y$ in $\mathcal{C}[\mathcal{W}^{-1}]$
should include not only morphisms in $\mathcal{C}$ but also ``formal
inverses'' of $\mathcal{W}$-arrows. The zigzag description handles
both layers at once.

\textbf{Reader:} And the universal property pins it down uniquely?

\textbf{Author:} Up to canonical isomorphism, yes. Two categories with
the same universal property are canonically isomorphic; this is the
abstract-nonsense way of saying the construction is well-defined.
\end{dialog}

\subsection{Formalizing}

\begin{definition}[Category with weak equivalences]
\label{def:cat-with-we}
A \term{category with weak equivalences} is a pair $(\mathcal{C}, \mathcal{W})$
where $\mathcal{C}$ is a category and $\mathcal{W} \subseteq
\mathrm{Mor}(\mathcal{C})$ is a subclass closed under composition and
containing all identities.

We typically also assume the \term{$2$-out-of-$3$ property}: if any two
of $f, g, gf$ are in $\mathcal{W}$, then so is the third.
\end{definition}

\begin{howtosay}
$\mathcal{W}$ \sayit{``the weak equivalences''}; the $2$-out-of-$3$
property \sayit{``two-out-of-three''} reflects that knowing two of the
three pieces of a triangle suffices to deduce the third. Sometimes the
stronger \term{$2$-out-of-$6$ property} (closure under sub-composition:
if $hg, gf$ are weak equivalences, so are $f, g, h, hgf$) is assumed.
\end{howtosay}

\begin{nameorigin}
The \term{$2$-out-of-$3$ property} is named for the combinatorial fact
it expresses: among three morphisms in a composable triple $f, g, gf$,
knowing two of them are in $\mathcal{W}$ forces the third. The
stronger \term{$2$-out-of-$6$ property} expresses a similar pattern
about six morphisms (two long composites and four atoms); it is
necessary for $\mathcal{W}$ to be \emph{saturated} — exactly equal to
the class of morphisms made invertible by localization.
\end{nameorigin}

\begin{example_thm}[Three running examples]
\label{ex:we-examples}
\textbf{(a) Topological spaces.} $\mathcal{W}_{\mathrm{Top}}$ = weak
homotopy equivalences: maps inducing isomorphisms on $\pi_n$ for all
$n \geq 0$. The category $\mathbf{Top}$ together with
$\mathcal{W}_{\mathrm{Top}}$ is a category with weak equivalences
satisfying $2$-out-of-$3$.

\textbf{(b) Chain complexes.} $\mathcal{W}_{\mathrm{Ch}(\mathcal{A})}$ =
\term{quasi-isomorphisms}: chain maps inducing isomorphisms on all
homology groups. Forms the entry point to homological algebra.

\textbf{(c) Simplicial sets.} $\mathcal{W}_{\mathrm{sSet}}$ = weak
equivalences: maps whose realization is a weak homotopy equivalence of
spaces. Chapter~41 will give a combinatorial characterization not
involving topological spaces.
\end{example_thm}

\begin{nameorigin}
A \term{quasi-isomorphism} is, literally, an ``almost-isomorphism'': a
chain map inducing isomorphisms on homology, but not necessarily
invertible at the chain level. ``Quasi-'' is from Latin \emph{quasi}
(``as if, almost''). The term was standardized in the 1950s--60s
homological algebra literature (Cartan--Eilenberg, Grothendieck).
\end{nameorigin}

\begin{readernote}
The three examples are sometimes called the ``trinity'' of homotopical
algebra: spaces, chain complexes, simplicial sets. Each has its own
notion of weak equivalence, but all three are connected by adjunctions
that turn out to be Quillen equivalences (Chapter~42), making their
homotopy categories ``the same.'' The localization construction below
handles all three uniformly.
\end{readernote}


% ═══════════════════════════════════════════════════════════════════════════
\section{The Gabriel--Zisman Localization}
\label{sec:gz-localization}
% ═══════════════════════════════════════════════════════════════════════════

\subsection{Formalizing}

\begin{nameorigin}
The \term{Gabriel--Zisman localization} is named for \emph{Pierre Gabriel}
(b.~1933, German--French mathematician) and \emph{Michel Zisman}
(b.~1932, French mathematician), whose 1967 book \emph{Calculus of
Fractions and Homotopy Theory} systematized the construction below in
the simplicial setting. The construction itself is older — versions
appear in ring theory, in Grothendieck's work on derived categories,
and in much earlier formal-inverse arguments — but Gabriel--Zisman
gave the definitive abstract treatment.
\end{nameorigin}

\begin{definition}[Localization]
\label{def:localization}
The \term{localization} of $(\mathcal{C}, \mathcal{W})$ is a category
$\mathcal{C}[\mathcal{W}^{-1}]$ together with a functor
\begin{equation*}
  Q : \mathcal{C} \to \mathcal{C}[\mathcal{W}^{-1}]
\end{equation*}
such that:
\begin{enumerate}
  \item $Q(w)$ is an isomorphism for every $w \in \mathcal{W}$.
  \item For any functor $F : \mathcal{C} \to \mathcal{D}$ that sends
        every $w \in \mathcal{W}$ to an isomorphism in $\mathcal{D}$,
        there is a unique $\bar F : \mathcal{C}[\mathcal{W}^{-1}] \to
        \mathcal{D}$ with $\bar F \circ Q = F$.
\end{enumerate}
The universal property determines $(\mathcal{C}[\mathcal{W}^{-1}], Q)$
uniquely up to canonical isomorphism.
\end{definition}

\begin{howtosay}
$Q$ \sayit{``the localization functor'' or ``the quotient''}; the
universal property \sayit{``functors out of $\mathcal{C}[\mathcal{W}^{-1}]$
correspond bijectively to functors out of $\mathcal{C}$ that already
invert $\mathcal{W}$''}. Some texts write $L_\mathcal{W}$ or just
$\gamma$ for the localization functor.
\end{howtosay}

\begin{theorem}[Gabriel--Zisman construction]
\label{thm:gz-construction}
For any $(\mathcal{C}, \mathcal{W})$, the localization
$\mathcal{C}[\mathcal{W}^{-1}]$ exists. It can be constructed as the
quotient of the free category on the graph $\mathcal{C} \sqcup
\mathcal{W}^{-1}$ — where $\mathcal{W}^{-1}$ adds a formal inverse edge
$w^{-1} : Y \to X$ for each $w : X \to Y$ in $\mathcal{W}$ — by the
relations:
\begin{itemize}
  \item Composition relations of $\mathcal{C}$;
  \item $w^{-1} \circ w = \mathrm{id}$ and $w \circ w^{-1} = \mathrm{id}$
        for each $w \in \mathcal{W}$.
\end{itemize}
\end{theorem}

\begin{nameorigin}
The morphisms of $\mathcal{C}[\mathcal{W}^{-1}]$ are called \term{zigzags}
because they look like zigzag paths: arrows alternately pointing forward
(in $\mathcal{C}$) and backward (a formal $w^{-1}$). The geometric
description is descriptive and goes back to the 1950s combinatorial
literature on free groups.
\end{nameorigin}

\begin{proof}[Proof sketch]
Construct $\mathcal{C}[\mathcal{W}^{-1}]$ as described. Morphisms are
\term{zigzags} of arrows
\begin{equation*}
  X \xrightarrow{f_1} Y_1 \xleftarrow{w_1} Z_1 \xrightarrow{f_2} \cdots \xleftarrow{w_n} Z_n \xrightarrow{f_{n+1}} Y,
\end{equation*}
with each backward arrow in $\mathcal{W}$, modulo the relations of the
free quotient. The functor $Q : \mathcal{C} \to \mathcal{C}[\mathcal{W}^{-1}]$
sends each morphism to itself (a length-one zigzag). The universal
property holds by direct verification: any $F$ inverting $\mathcal{W}$
extends uniquely to zigzags by sending each $w^{-1}$ to $F(w)^{-1}$.
\qed
\end{proof}

\begin{remark}[Size issues]
\label{rem:size-issues}
The free quotient construction is set-theoretically problematic when
$\mathcal{C}$ is large: hom-classes of $\mathcal{C}[\mathcal{W}^{-1}]$
may not be small even when $\mathcal{C}$'s are. The model-categorical
machinery of Chapter~41 will resolve this for the examples we care
about by producing \emph{explicit small-hom-set models} for
$\mathcal{C}[\mathcal{W}^{-1}]$ (the homotopy category as morphisms
between bifibrant objects).
\end{remark}

\begin{readernote}
The Gabriel--Zisman construction proves existence but is rarely useful
for \emph{computation}: a zigzag of length 17 may represent the same
morphism as a zigzag of length 3, but checking equivalence requires
following a long chain of relations. The calculus of fractions
(\S\ref{sec:calculus-of-fractions}) is the standard remedy: under
suitable conditions, every zigzag reduces to length 2, giving a
concrete description of $\mathcal{C}[\mathcal{W}^{-1}]$'s morphisms.
\end{readernote}


% ═══════════════════════════════════════════════════════════════════════════
\section{Calculus of Fractions}
\label{sec:calculus-of-fractions}
% ═══════════════════════════════════════════════════════════════════════════

\subsection{Informally}

The Gabriel--Zisman construction describes a morphism of
$\mathcal{C}[\mathcal{W}^{-1}]$ as a zigzag of arbitrary length. Under
additional hypotheses on $\mathcal{W}$ — the \term{calculus of fractions}
— every zigzag reduces to a single span (or cospan), yielding a much
more manageable description.

\begin{nameorigin}
The \term{calculus of fractions} is named by analogy with the
calculus of fractions in (non-commutative) ring theory: \emph{Øystein
Ore} (1899--1968, Norwegian mathematician) gave conditions in 1931
under which a non-commutative ring $R$ has a ring of fractions
$S^{-1} R$. The \emph{Ore condition} below is the same combinatorial
condition lifted to categories. ``Fractions'' because the morphisms of
the localization look like $w^{-1} f$ (formal ``$f$ over $w$''), like
ratios in algebra.
\end{nameorigin}

\subsection{Formalizing}

\begin{definition}[Left calculus of fractions]
\label{def:left-fractions}
A class $\mathcal{W}$ admits a \term{left calculus of fractions} if:
\begin{enumerate}
  \item \emph{Identities and composition.} $\mathcal{W}$ contains all
        identities and is closed under composition.
  \item \emph{Ore condition.} For every span $X \xleftarrow{w} A
        \xrightarrow{f} B$ with $w \in \mathcal{W}$, there exists a
        commutative square
        \begin{equation*}
          \begin{tikzcd}[column sep=large]
            A \arrow[r, "f"] \arrow[d, "w"'] & B \arrow[d, "w'"] \\
            X \arrow[r, "f'"'] & Y
          \end{tikzcd}
        \end{equation*}
        with $w' \in \mathcal{W}$.
  \item \emph{Cancellation.} If $fw = gw$ with $w \in \mathcal{W}$,
        there exists $w' \in \mathcal{W}$ such that $w'f = w'g$.
\end{enumerate}
\end{definition}

\begin{howtosay}
``Calculus of fractions'' \sayit{``calculus of fractions''}; the
\emph{Ore condition} \sayit{``the Ore (or ``or-AY'') condition''},
sometimes \sayit{``the left Ore condition''} to distinguish from its
dual. Sometimes also called \term{Gabriel's axiom} or
\term{Hartshorne's localization theorem} depending on context.
\end{howtosay}

\begin{remark}[Right calculus]
The dual notion of \term{right calculus of fractions} reverses arrows:
the Ore condition completes \emph{cospans} $X \xrightarrow{f} B
\xleftarrow{w} A$ to commutative squares. Each side of the calculus
gives a different reduction (left fractions are spans $X \xleftarrow{w}
\cdot \xrightarrow{f} Y$; right fractions are cospans).
\end{remark}

\begin{theorem}[Reduction to fractions]
\label{thm:fractions-reduction}
If $\mathcal{W}$ admits a left calculus of fractions, every morphism of
$\mathcal{C}[\mathcal{W}^{-1}]$ can be represented by a single
\term{left fraction}
\begin{equation*}
  X \xleftarrow{w} A \xrightarrow{f} Y, \qquad w \in \mathcal{W},
\end{equation*}
and two left fractions $(w, f)$ and $(w', f')$ represent the same
morphism iff there is a common refinement
\begin{equation*}
  \begin{tikzcd}[column sep=small, row sep=small]
    & A \arrow[ldd, "w"', bend right] \arrow[rdd, "f", bend left] & \\
    & B \arrow[d, "u"'] \arrow[u, "v"', dashed] & \\
    X & A' \arrow[l, "w'"] \arrow[r, "f'"'] & Y
  \end{tikzcd}
\end{equation*}
with $u, v \in \mathcal{W}$ commuting both squares.
\end{theorem}

\begin{proof}[Proof sketch]
Every zigzag $f_1, w_1^{-1}, f_2, w_2^{-1}, \ldots$ is reduced by
repeatedly applying the Ore condition (interchange a backward $w^{-1}$
with the following forward $f$ by completing the span). Zigzag length
strictly decreases at each step; the procedure terminates with a single
left fraction. Equivalence of two fractions is the equivalence relation
generated by common refinement; the cancellation axiom guarantees that
the relation is transitive.
\qed
\end{proof}

\begin{remark}[Composition of fractions]
With the calculus in hand, composition is concrete: given fractions $X
\xleftarrow{w} A \xrightarrow{f} Y$ and $Y \xleftarrow{v} B
\xrightarrow{g} Z$, complete the span $A \xrightarrow{f} Y \xleftarrow{v}
B$ to a commutative square by Ore, obtaining $A \xleftarrow{v'} C
\xrightarrow{f'} B$ with $v' \in \mathcal{W}$; the composite is then
$X \xleftarrow{wv'} C \xrightarrow{gf'} Z$.
\end{remark}

\begin{readernote}
The Ore condition is best remembered geometrically: it says \emph{spans
with one $\mathcal{W}$-leg can be completed to commutative squares with
$\mathcal{W}$-legs on both top and bottom}. This is exactly the
``rewriting rule'' used to push backwards arrows to the left of
forwards arrows, simplifying zigzags. Without the Ore condition, the
rewriting gets stuck and zigzags remain arbitrarily long.
\end{readernote}


% ═══════════════════════════════════════════════════════════════════════════
\section{Examples}
\label{sec:localization-examples}
% ═══════════════════════════════════════════════════════════════════════════

\subsection{The derived category of an abelian category}

\begin{example_thm}[Derived category $\mathcal{D}(\mathcal{A})$]
\label{ex:derived-cat}
Let $\mathcal{A}$ be an abelian category and $\mathbf{Ch}(\mathcal{A})$
its category of chain complexes. The \term{derived category}
\begin{equation*}
  \mathcal{D}(\mathcal{A}) \;=\; \mathbf{Ch}(\mathcal{A})\bigl[\mathcal{W}^{-1}\bigr]
\end{equation*}
is the localization at $\mathcal{W} =$ quasi-isomorphisms.

The intermediate \emph{homotopy category} $\mathbf{K}(\mathcal{A}) =
\mathbf{Ch}(\mathcal{A})/\sim$ (chain maps modulo chain homotopy) admits
a calculus of fractions with respect to quasi-isomorphisms — both left
and right. Therefore $\mathcal{D}(\mathcal{A}) =
\mathbf{K}(\mathcal{A})[\mathcal{W}^{-1}]$ has the concrete description:
morphisms are spans (or cospans) of chain maps with quasi-isomorphism
legs, modulo refinement.
\end{example_thm}

\begin{nameorigin}
The \term{derived category} $\mathcal{D}(\mathcal{A})$ was introduced
by Grothendieck's student \emph{Jean-Louis Verdier} (1935--1989) in his
1967 thesis. ``Derived'' because morphisms in $\mathcal{D}(\mathcal{A})$
are equivalent to \emph{derived functor} information (Tor, Ext, etc.):
in $\mathcal{D}(\mathcal{A})$, the natural map $\mathrm{Hom}(M, N) \to
\mathrm{Hom}(M^\bullet, N^\bullet)$ encodes Tor/Ext via the chain
representatives.
\end{nameorigin}

\subsection{The homotopy category of spaces}

\begin{example_thm}[$\mathrm{Ho}(\mathbf{Top})$]
\label{ex:ho-top}
The homotopy category $\mathrm{Ho}(\mathbf{Top}) = \mathbf{Top}[\mathcal{W}^{-1}]$
inverts weak homotopy equivalences. The calculus of fractions does
\emph{not} apply directly in $\mathbf{Top}$; constructing the
localization requires explicit replacement of spaces by CW complexes
(an essentially model-categorical operation). After replacement, every
morphism of $\mathrm{Ho}(\mathbf{Top})$ between CW complexes is a
homotopy class of continuous maps, by Whitehead's theorem (Chapter~41).
\end{example_thm}

\begin{nameorigin}
\term{Whitehead's theorem} is named for \emph{J.H.C.\ Whitehead}
(1904--1960, British algebraic topologist; not to be confused with
Alfred North Whitehead). Whitehead's theorem says that a weak homotopy
equivalence between CW complexes is an honest homotopy equivalence —
a non-trivial fact tying the weak (homotopy-group-iso) and strong
(homotopy-inverse) notions together for sufficiently nice spaces.
\end{nameorigin}

\subsection{Quillen's localization}

\begin{remark}[Model-categorical localization]
A \term{model category} (Chapter~41) is a category with three classes
of maps — weak equivalences, cofibrations, fibrations — satisfying
axioms that together provide an explicit construction of
$\mathcal{C}[\mathcal{W}^{-1}]$: the homotopy category
$\mathrm{Ho}(\mathcal{C})$ has the same objects as $\mathcal{C}$, and
its morphisms between bifibrant objects $X$ (cofibrant and fibrant) are
homotopy classes $[X, Y]/\sim$. The model-categorical localization is
small-hom-class by construction, sidestepping the size issue of
Remark~\ref{rem:size-issues}.
\end{remark}

\begin{readernote}
The pattern across all three examples — chain complexes, topological
spaces, simplicial sets — is the same: a raw category with weak
equivalences is too rigid to localize directly, but admits a
\emph{replacement} (chain homotopy classes, CW complexes, bifibrant
simplicial sets) where the calculus of fractions does work, or
equivalently, where homotopy classes give an honest description of
localized morphisms. Chapter~41 axiomatizes this pattern.
\end{readernote}


% ═══════════════════════════════════════════════════════════════════════════
\section{Exercises}
\label{sec:localization-exercises}
% ═══════════════════════════════════════════════════════════════════════════

\begin{exercisesbox}
\begin{qa}
\qaitem{Define a category with weak equivalences.}{A pair $(\mathcal{C}, \mathcal{W})$ where $\mathcal{W} \subseteq \mathrm{Mor}(\mathcal{C})$ contains identities and is closed under composition.}
\qaitem{What is the $2$-out-of-$3$ property?}{If two of $f, g, gf$ are in $\mathcal{W}$, the third is too.}
\qaitem{Where does ``localization'' come from?}{Algebra: localizing a ring $R$ at a multiplicative set $S$ inverts the elements of $S$. ``Local'' because the result captures behavior near a prime ideal. Lifted to categories in the 1960s--70s.}
\qaitem{Where does ``weak equivalence'' come from?}{Contrast with strict isomorphism: a weak equivalence is iso ``up to homotopy or some other higher equivalence,'' not on the nose. Standard in algebraic topology.}
\qaitem{Define the localization $\mathcal{C}[\mathcal{W}^{-1}]$ by its universal property.}{A category with a functor $Q : \mathcal{C} \to \mathcal{C}[\mathcal{W}^{-1}]$ inverting $\mathcal{W}$, universal: any functor $F$ inverting $\mathcal{W}$ factors uniquely through $Q$.}
\qaitem{Why does the localization always exist?}{Gabriel--Zisman: build the free category on $\mathcal{C} \sqcup \mathcal{W}^{-1}$ and quotient by the appropriate relations.}
\qaitem{Who are Gabriel and Zisman?}{Pierre Gabriel (b. 1933) and Michel Zisman (b. 1932), authors of the 1967 book \emph{Calculus of Fractions and Homotopy Theory}, which systematized the localization construction.}
\qaitem{What is a zigzag in $\mathcal{C}[\mathcal{W}^{-1}]$?}{A sequence of forward arrows in $\mathcal{C}$ and backward arrows in $\mathcal{W}^{-1}$ from $X$ to $Y$ — the generic morphism description from the Gabriel--Zisman construction.}
\qaitem{Why is the localization set-theoretically subtle?}{Hom-classes in $\mathcal{C}[\mathcal{W}^{-1}]$ may not be small (proper classes of zigzag equivalence classes) even when those of $\mathcal{C}$ are.}
\qaitem{Define a left calculus of fractions.}{$\mathcal{W}$ contains identities and closes under composition; the Ore condition completes spans with $\mathcal{W}$-leg to commutative squares with $\mathcal{W}$-leg; cancellation: $fw = gw$ implies $w'f = w'g$ for some $w' \in \mathcal{W}$.}
\qaitem{Who is Ore?}{Øystein Ore (1899--1968), Norwegian mathematician. Gave conditions in 1931 for non-commutative ring localization; the Ore condition above is the categorical version.}
\qaitem{What does the Ore condition replace zigzags with?}{Single left fractions $X \xleftarrow{w} A \xrightarrow{f} Y$ with $w \in \mathcal{W}$. Every zigzag reduces to a span by repeated application of Ore.}
\qaitem{When are two left fractions equivalent?}{When they admit a common refinement: a span $A' \to A$ and $A' \to B$ both in $\mathcal{W}$ that fits between the two.}
\qaitem{How do you compose left fractions?}{Given $X \xleftarrow{w} A \xrightarrow{f} Y$ and $Y \xleftarrow{v} B \xrightarrow{g} Z$, Ore-complete $f, v$ to a square with $v', f'$, obtaining $X \xleftarrow{wv'} C \xrightarrow{gf'} Z$.}
\qaitem{What is the derived category of an abelian category?}{$\mathcal{D}(\mathcal{A}) = \mathbf{Ch}(\mathcal{A})[\mathcal{W}^{-1}]$ where $\mathcal{W}$ = quasi-isomorphisms (maps inducing isomorphism on homology).}
\qaitem{Who introduced derived categories?}{Jean-Louis Verdier (Grothendieck's student) in his 1967 thesis. ``Derived'' because morphisms encode derived-functor information (Tor, Ext).}
\qaitem{What is a quasi-isomorphism?}{A chain map inducing isomorphisms on homology. ``Quasi-'' (Latin \emph{quasi}, ``almost'') because it's almost an isomorphism — invertible at the level of homology, not necessarily at the chain level.}
\qaitem{Why is the chain-homotopy category $\mathbf{K}(\mathcal{A})$ a useful stepping stone?}{Because quasi-isomorphisms admit a calculus of fractions in $\mathbf{K}(\mathcal{A})$ (but not in $\mathbf{Ch}(\mathcal{A})$). So $\mathcal{D}(\mathcal{A}) = \mathbf{K}(\mathcal{A})[\mathcal{W}^{-1}]$ has the concrete left/right-fractions description.}
\qaitem{What is $\mathrm{Ho}(\mathbf{Top})$?}{$\mathbf{Top}[\mathcal{W}^{-1}]$ with $\mathcal{W}$ = weak homotopy equivalences. Concretely, the homotopy category of CW complexes (by Whitehead).}
\qaitem{Who is Whitehead?}{J.H.C.\ Whitehead (1904--1960), British algebraic topologist. His theorem says weak homotopy equivalences between CW complexes are honest homotopy equivalences — tying weak and strong notions.}
\qaitem{Why isn't the calculus of fractions enough for $\mathbf{Top}$?}{The Ore condition fails: completing a span involving a weak homotopy equivalence may require a CW replacement that lives outside $\mathbf{Top}$ in any natural sense. Model categories handle this systematically.}
\qaitem{What's the model-category alternative to Gabriel--Zisman?}{Replace objects by bifibrant ones and define $\mathrm{Ho}(\mathcal{C})(X, Y) = [X, Y]/\sim$ (homotopy classes). Gives the same localization but with small hom-classes by construction.}
\qaitem{State the universal property of $Q$ in equation form.}{$\mathrm{Fun}(\mathcal{C}[\mathcal{W}^{-1}], \mathcal{D}) \xrightarrow{\, -\circ Q\,} \mathrm{Fun}_{\mathcal{W}\text{-inv}}(\mathcal{C}, \mathcal{D})$ is an isomorphism of categories.}
\qaitem{What does inverting $\mathcal{W}$ feel like categorically?}{The functor $Q$ is essentially surjective but not faithful; it identifies parallel morphisms that differ by a chain of $\mathcal{W}$-zigzags.}
\qaitem{Does $\mathcal{C}[\mathcal{W}^{-1}]$ have the same objects as $\mathcal{C}$?}{Yes, by construction. The localization changes only morphisms.}
\qaitem{What's an example where the calculus of fractions fails?}{$\mathbf{Top}$ with weak homotopy equivalences: the Ore square cannot in general be filled within $\mathbf{Top}$ without CW replacement.}
\qaitem{Why is the $2$-out-of-$6$ property sometimes assumed?}{It makes a class of weak equivalences ``saturated'': closed under recovering weak equivalences from longer composites. Implies that $\mathcal{W}$ contains exactly the morphisms made invertible by $Q$.}
\qaitem{Is every isomorphism in $\mathcal{C}$ in $\mathcal{W}$?}{Not required by the bare definition, but yes if $\mathcal{W}$ is \emph{saturated} (closed under invertible-becoming-after-localization). Most natural classes of weak equivalences contain all isomorphisms.}
\qaitem{What's the relationship between $\mathcal{C}[\mathcal{W}^{-1}]$ and Chapter~43's quasi-categories?}{$\mathcal{C}[\mathcal{W}^{-1}]$ is the homotopy 1-category of an underlying $\infty$-category — the latter retains higher coherence data that the localization discards. Volume IV's $\infty$-categorical Yoneda (Chapter~48) lives at this richer level.}
\qaitem{Why is localization a left adjoint?}{$Q : \mathcal{C} \to \mathcal{C}[\mathcal{W}^{-1}]$ has the universal property of an initial object in the comma category of $\mathcal{W}$-inverting functors out of $\mathcal{C}$; this is left-adjoint-style universality on the level of the $2$-category $\mathbf{Cat}$.}
\end{qa}
\end{exercisesbox}


% ═══════════════════════════════════════════════════════════════════════════
\section*{Chapter Summary}
\addcontentsline{toc}{section}{Chapter Summary}
% ═══════════════════════════════════════════════════════════════════════════

\begin{summarybox}
\textbf{Category with weak equivalences.}\quad $(\mathcal{C}, \mathcal{W})$
with $\mathcal{W}$ closed under composition and containing identities,
typically with the $2$-out-of-$3$ property. Origin: ``localization''
from algebra/geometry; ``weak equivalence'' from algebraic topology.

\textbf{Localization (Gabriel--Zisman, 1967).}\quad
$\mathcal{C}[\mathcal{W}^{-1}]$ is universal among categories receiving
a $\mathcal{W}$-inverting functor from $\mathcal{C}$. Always exists,
with morphisms presented as zigzags modulo relations.

\textbf{Calculus of fractions (Ore, 1931 + Gabriel--Zisman).}\quad
An Ore-type condition under which zigzags reduce to single left (or
right) fractions. Composition becomes concrete; equivalence becomes
common refinement.

\textbf{Three examples.}\quad The derived category $\mathcal{D}(\mathcal{A})$
(Verdier 1967); the homotopy category $\mathrm{Ho}(\mathbf{Top})$
(Whitehead 1949); Quillen's model categories (1967). The last unifies
the previous by giving an explicit small-hom-class construction via
bifibrant replacement (Chapter~41).
\end{summarybox}


% ═══════════════════════════════════════════════════════════════════════════
\section*{Formal Restatement}
\addcontentsline{toc}{section}{Formal Restatement}
% ═══════════════════════════════════════════════════════════════════════════

\begin{formalrestatement}
\textbf{Localization.}\quad For $(\mathcal{C}, \mathcal{W})$, the
localization $(Q, \mathcal{C}[\mathcal{W}^{-1}])$ satisfies:
\begin{equation*}
  \mathrm{Fun}_{\mathcal{W}\text{-inv}}(\mathcal{C}, \mathcal{D}) \;\cong\;
  \mathrm{Fun}(\mathcal{C}[\mathcal{W}^{-1}], \mathcal{D}),
  \qquad F \longmapsto F \circ Q.
\end{equation*}

\medskip
\textbf{Left calculus of fractions.}\quad $\mathcal{W}$ admits left
fractions if:
\begin{itemize}
  \item $\mathcal{W}$ has identities and closes under composition.
  \item (Ore) Every span $X \xleftarrow{w} A \xrightarrow{f} B$ with
        $w \in \mathcal{W}$ completes to a commutative square with the
        right leg in $\mathcal{W}$.
  \item (Cancellation) $fw = gw$ with $w \in \mathcal{W}$ implies
        $w'f = w'g$ for some $w' \in \mathcal{W}$.
\end{itemize}

\medskip
\textbf{Fractions theorem.}\quad If $\mathcal{W}$ admits left fractions,
$\mathcal{C}[\mathcal{W}^{-1}](X, Y)$ is the set of equivalence classes
of left fractions $X \xleftarrow{w} A \xrightarrow{f} Y$ under common
refinement.

\medskip
\textbf{Three running examples.}\quad
\begin{itemize}
  \item $\mathbf{Ch}(\mathcal{A})$ with quasi-isomorphisms: $\mathcal{D}(\mathcal{A})$.
  \item $\mathbf{Top}$ with weak homotopy equivalences: $\mathrm{Ho}(\mathbf{Top})$.
  \item $\mathbf{sSet}$ with realization-weak-equivalences: $\mathrm{Ho}(\mathbf{sSet})$.
\end{itemize}
\end{formalrestatement}
